\lektionstrennseite{01}{01}{Sūrat al-Fātiḥa und die kurzen Sūren}

\section*{Lektion 1: Sūrat al-Fātiḥa und die kurzen Sūren}

\begin{lernziele}
\begin{itemize}[leftmargin=1.4em]
  \item erklären können, welche Sūren in dieser Lektion gemeint sind und warum gerade mit
        ihnen begonnen wird,
  \item Sūrat al-Fātiḥa und die kurzen Sūren (von Sūrat az-Zalzala bis Sūrat an-Nās) korrekt
        rezitieren und nach und nach auswendig lernen,
  \item benennen können, mit welchem Tafsīr-Werk ein Anfänger sinnvollerweise beginnt und
        warum,
  \item die drei Grundhaltungen gegenüber dem Qur’ān unterscheiden können und wissen,
        welche davon der richtige Weg ist.
\end{itemize}
\end{lernziele}

\begin{shaykhsagt}
{\arabicfont\begin{Arabic}
سورة الفاتحة و ما أمكن من قصار السُّور، من سورة الزلزلة إلى سورة الناس، تلقيناً، وتصحيحاً
للقراءة، وتحفيظاً، و شرحاً لما يجب فهمه.
\end{Arabic}}

\vspace{2mm}
\emph{Sūrat al-Fātiḥa und, soweit wie möglich, auch die kurzen Sūren – von Sūrat az-Zalzala bis
Sūrat an-Nās – werden gelernt; die Rezitation wird korrigiert, sie werden auswendig gelernt, und
es wird erläutert, was davon verstanden werden muss.}
\end{shaykhsagt}

\subsection*{Erläuterung}

Der Shaykh \ar{رحمه الله} beginnt sein Werk bewusst nicht mit einer trockenen Definition, sondern
mit dem Qur’ān selbst. Bevor irgendetwas anderes gelernt wird, soll ein Muslim zuerst
Sūrat al-Fātiḥa und die kurzen, oft rezitierten Sūren aus dem letzten Teil des Qur’ān beherrschen
– vom Sinn her verstehen, richtig aussprechen und auswendig können. Das ist keine Nebensache:
Diese Sūren werden in jedem Gebet rezitiert, und ohne sie ist kein Gebet vollständig.

\subsubsection*{Wie lernt man den Qur’ān richtig?}

Wie schon die frühen Muslime es taten, sollte man sich angewöhnen, regelmäßig – zum Beispiel
täglich etwa zehn Verse – auswendig zu lernen. Dabei liest man parallel dazu eine kurze,
verständliche Tafsīr (Koranauslegung) der jeweiligen Verse und bittet Allah \ar{ﷻ} darum, einem
dabei zu helfen, das Gelernte auch umzusetzen.

\subsubsection*{Mit welchem Tafsīr-Werk sollte man beginnen?}

Einem Anfänger wird empfohlen, mit dem Tafsīr-Werk von Shaykh ʿAbd ar-Raḥmān Ibn Nāṣir
as-Saʿdī \ar{رحمه الله} zu beginnen.

\begin{merke}
Vier Gründe, warum gerade dieses Werk für den Einstieg empfohlen wird:
\begin{enumerate}[leftmargin=1.4em]
  \item Die Gelehrten empfahlen es und beschäftigten sich besonders damit.
  \item Es ist kurz und dadurch für Anfänger gut lesbar.
  \item Seine Sprache ist einfach, klar und nicht mehrdeutig.
  \item Der Verfasser \ar{رحمه الله} legt darin besonderen Wert auf den Tawḥīd, die Einzigkeit
        Allahs.
\end{enumerate}
\end{merke}

\subsubsection*{Wie verhalten sich Menschen gegenüber dem Qur’ān?}

Im Umgang mit dem Qur’ān lassen sich Menschen grob in drei Gruppen einteilen: zwei
Übertreibungen an den jeweiligen Enden und einen richtigen Mittelweg dazwischen.

\begin{center}
\renewcommand{\arraystretch}{1.5}
\begin{tabularx}{\linewidth}{@{}>{\raggedright\arraybackslash}X >{\raggedright\arraybackslash}X >{\raggedright\arraybackslash}X@{}}
\toprule
\textbf{\color{DarkGreen}Vernachlässigung} & \textbf{\color{Accent}Der Mittelweg (richtig)} & \textbf{\color{DarkGreen}Übertreibung} \\
\midrule
Der Qur’ān wird vernachlässigt – beim Lesen, beim Auswendiglernen, beim Nachdenken darüber,
bei der Umsetzung im Alltag oder sogar dabei, sich mit seinen Versen heilen zu lassen.
&
Der Qur’ān wird regelmäßig rezitiert und auswendig gelernt; man bittet Allah \ar{ﷻ} darum,
beim Umsetzen des Gelernten zu helfen.
&
Der Qur’ān wird zwar rezitiert und auswendig gelernt, aber ohne über seine Bedeutung
nachzudenken und ohne ihn im eigenen Leben umzusetzen.
\\
\bottomrule
\end{tabularx}
\end{center}

Allah \ar{ﷻ} beschreibt die erste Haltung mit den Worten des Gesandten \ar{ﷺ} am Jüngsten Tag:

\begin{quranvers}{Sūrah al-Furqān, 25:30}
\emph{„Und der Gesandte sagt: ‚O mein Herr, mein Volk hat diesen Qur’ān verlassen und
gemieden.‘“}
\end{quranvers}

Der Prophet \ar{ﷺ} warnte zudem eindringlich vor der zweiten Haltung – vor Menschen, die den
Qur’ān zwar rezitieren, ihre Kehlen aber nicht wirklich erreicht, das heißt: ihr Herz nicht wirklich
berührt. Sie tragen den Qur’ān auf der Zunge, ohne dass er sich in ihrem Handeln zeigt.

\begin{zusammenfassung}
\begin{itemize}[leftmargin=1.4em]
  \item Der Shaykh \ar{رحمه الله} beginnt bewusst mit dem Qur’ān: Sūrat al-Fātiḥa und die kurzen
        Sūren von az-Zalzala bis an-Nās sollen gehört, richtig rezitiert, auswendig gelernt und
        inhaltlich verstanden werden.
  \item Empfohlene Lernmethode: regelmäßig (z.\,B. täglich) einen kleinen Abschnitt auswendig
        lernen und dabei eine kurze, verständliche Tafsīr lesen.
  \item Für den Einstieg wird das Tafsīr-Werk von Shaykh as-Saʿdī \ar{رحمه الله} empfohlen: kurz,
        klar und mit besonderem Fokus auf den Tawḥīd.
  \item Der richtige Umgang mit dem Qur’ān liegt in der Mitte zwischen zwei Extremen:
        Vernachlässigung auf der einen und bloßes Rezitieren ohne Nachdenken und Umsetzung
        auf der anderen Seite.
\end{itemize}
\end{zusammenfassung}
